% Wersja źródłowa: audyt II, jednostka 21

\section{Splot dystrybucji i regularyzacja}
\label{sec:splot21}

\providecommand{\atlasref}[2]{#2}


W całym rozdziale $d\ge1$ jest liczbą całkowitą,
$\K\in\{\R,\C\}$, a dystrybucje
są ciągłymi $\K$-liniowymi funkcjonałami na
$\mathcal D(\R^d;\K)$. Rozpatrujemy splot na całej
$\R^d$; dla ogólnego zbioru otwartego $\Omega$ nie ma
jednej operacji splotu dwóch dystrybucji z
$\mathcal D'(\Omega)$ bez dodatkowych ustaleń.
Oznaczymy przez $\mathcal E'(\R^d)$ przestrzeń dystrybucji
o zwartym nośniku. Nie należy mylić jej ze zbiorem funkcji
testowych $\mathcal D(\R^d)$. Pominięcie $\K$ w oznaczeniu
przestrzeni oznacza pracę nad ustalonym ciałem $\K$.
Dla $f\in L^1_{\mathrm{loc}}(\R^d;\K)$ piszemy $T_f$ na
odpowiadającą mu dystrybucję regularną. W zapisie splotu
$T*\rho$ dla $\rho\in\mathcal D(\R^d)$ rozumiemy
$T*T_\rho$; jeżeli splot jest regularny, jego funkcję
reprezentującą będziemy wskazywać jawnie.
Dla zbiorów $K,F\subseteq\R^d$ oznaczamy
$K+F=\{k+f:k\in K,\ f\in F\}$ i
$K-F=\{k-f:k\in K,\ f\in F\}$.

\subsection{Iterowane działania dystrybucji i definicja splotu}

\begin{lemat}[Przestawianie dwóch działań dystrybucji]
\label{lem:iterated-distributions21}
Niech $k,l\ge1$ będą liczbami całkowitymi,
$T\in\mathcal D'(\R^k)$,
$S\in\mathcal D'(\R^l)$ i
$\Phi\in\mathcal D(\R^{k+l})$. Funkcje
\[
 x\longmapsto\langle S_y,\Phi(x,y)\rangle,
 \qquad
 y\longmapsto\langle T_x,\Phi(x,y)\rangle
\]
należą odpowiednio do $\mathcal D(\R^k)$ i
$\mathcal D(\R^l)$. Ponadto
\[
 \langle T_x,\langle S_y,\Phi(x,y)\rangle\rangle
 =\langle S_y,\langle T_x,\Phi(x,y)\rangle\rangle.
\]
\end{lemat}
\begin{proof}
Ustalmy zwarte prostokąty $K_x\subseteq\R^k$ i
$K_y\subseteq\R^l$, których wnętrza zawierają odpowiednie
rzuty $\supp\Phi$. Dla każdego $x$ funkcja $\Phi(x,\cdot)$
należy do $\mathcal D_{K_y}$. Ilorazy różnicowe względem
$x_j$ zbiegają w każdej seminormie $p_m$ tej przestrzeni
do $\partial_{x_j}\Phi(x,\cdot)$: po różniczkowaniu
względem $y$ błąd jest ograniczony przez moduł ciągłości
odpowiedniej pochodnej na iloczynie zwartego otoczenia $x$
i $K_y$, na mocy twierdzenia Newtona--Leibniza.
Powtarzając ten argument dla kolejnych pochodnych i
stosując ciągłość $S$ na $\mathcal D_{K_y}$, otrzymujemy
\[
 \partial_x^\alpha\langle S_y,\Phi(x,y)\rangle
 =\langle S_y,\partial_x^\alpha\Phi(x,y)\rangle.
\]
Ciągłość tych pochodnych wynika z jednostajnej ciągłości
odpowiednich pochodnych $\Phi$ na tych samych zbiorach.
Analogicznie postępujemy z $T$. Nośniki obu otrzymanych
funkcji mieszczą się w odpowiednich rzutach $\supp\Phi$,
więc są to funkcje testowe.

Istnieją $C_T,C_S>0$ i $m_T,m_S\in\N_0$ takie, że
$|\langle T,u\rangle|\le C_Tp_{m_T}(u)$ dla
$u\in\mathcal D_{K_x}$ i analogicznie dla $S$ na
$\mathcal D_{K_y}$. Z powyższego wzoru na pochodne mamy
\[
 |\langle T_x,\langle S_y,\Phi(x,y)\rangle\rangle|
 \le C_TC_S p_{m_T+m_S}(\Phi).
\]
To samo oszacowanie zachodzi dla odwrotnej kolejności.
W szczególności oba działania są ciągłe na każdym etapie
o ustalonym zwartym nośniku.

Pozostaje uzasadnić równość. Dla $\Phi(x,y)=u(x)v(y)$,
gdzie $u,v$ są funkcjami testowymi, obie strony są równe
$\langle T,u\rangle\langle S,v\rangle$.
Pokażemy, że skończone sumy takich funkcji aproksymują
dowolną $\Phi$ w jednym etapie przestrzeni $\mathcal D$.
Wybierzmy nieujemne rzeczywiste funkcje testowe
$\rho$ na $\R^k$ i $\sigma$ na $\R^l$, o całkach $1$
i nośnikach w kulach jednostkowych, przez normalizację
funkcji odcięcia z Lematu~\ref{lem:cutoffs17}.
Niech $\rho_\varepsilon(x)=\varepsilon^{-k}\rho(x/\varepsilon)$
i $\sigma_\varepsilon(y)=\varepsilon^{-l}\sigma(y/\varepsilon)$.
Dla $0<\varepsilon\le1$ określamy
\[
 \Phi_\varepsilon(x,y)=
 \int_{\R^k\times\R^l}\Phi(u,v)
 \rho_\varepsilon(x-u)\sigma_\varepsilon(y-v)
 \,\mathrm d\lambda_{k+l}(u,v).
\]
Wszystkie te funkcje mają nośniki w ustalonym zwartym
prostokącie $L_x\times L_y$, gdzie $L_x,L_y$ zawierają
odpowiednio $K_x+\overline{B(0,1)}$ i
$K_y+\overline{B(0,1)}$. Po zmianie zmiennych całka
jest średnią przesuniętych funkcji $\Phi$ z wagą
$\rho(u)\sigma(v)$. Dla każdego $m$ jednostajna ciągłość
pochodnych $\Phi$ rzędu co najwyżej $m$ daje
$p_m(\Phi_\varepsilon-\Phi)\to0$.

Dla ustalonego $\varepsilon$ całkę definiującą
$\Phi_\varepsilon$ można ograniczyć do $K_x\times K_y$.
Jej sumy Riemanna są skończonymi sumami funkcji postaci
$c\,\rho_\varepsilon(x-u)\sigma_\varepsilon(y-v)$.
Dla każdego skończonego $m$ sumy te zbiegają w $p_m$:
pochodne względem $(x,y)$ można przenieść pod całkę,
a funkcje podcałkowe po różniczkowaniu są jednostajnie
ciągłe na zwartym iloczynie prostokątów zmiennych
$(x,y)$ i $(u,v)$. Poza $L_x\times L_y$ zarówno sumy,
jak i całka znikają. Dla $\varepsilon=1/n$ wybieramy
sumę $F_n$ z błędem $p_n(F_n-\Phi_{1/n})<1/n$.
Wtedy dla każdego ustalonego $m$ i $n\ge m$ mamy
$p_m(F_n-\Phi)\le 1/n+p_m(\Phi_{1/n}-\Phi)\to0$.
Nośniki wszystkich $F_n$ zawierają się w jednym zwartym
zbiorze $L_x\times L_y$, a funkcje te są sumami
iloczynów funkcji testowych nad ustalonym ciałem $\K$.
Stosując uzyskane oszacowania, tym razem na etapie
o nośniku zawartym w $L_x\times L_y$, do różnicy $F_n-\Phi$,
przechodzimy do granicy w równości obu działań na $F_n$.
\end{proof}

\begin{lemat}[Operator testowy indukowany dystrybucją]
\label{lem:convolution-transpose21}
Niech $T\in\mathcal E'(\R^d)$, $K=\supp T$, i niech
$\varphi\in\mathcal D(\R^d)$. Funkcja
\[
 (A_T\varphi)(y)=\langle T_x,\varphi(x+y)\rangle
\]
należy do $\mathcal D(\R^d)$ i spełnia
$\supp(A_T\varphi)\subseteq\supp\varphi-K$.
Operator $A_T:\mathcal D(\R^d)\to\mathcal D(\R^d)$
jest liniowy i ciągły, a obrazy zbiorów ograniczonych
są ograniczone.
\end{lemat}
\begin{proof}
W działaniu $T_x$ korzystamy z przedłużenia z
Lematu~\ref{lem:compact-action21}. Dla $y$ z ustalonego
zbioru zwartego funkcje $x\mapsto\varphi(x+y)$ mają
wspólny zwarty nośnik; można więc różniczkować po $y$:
$\partial_y^\beta(A_T\varphi)(y)=
\langle T_x,\partial^\beta\varphi(x+y)\rangle$.
Jeżeli $y\notin\supp\varphi-K$, funkcja
$\varphi(\cdot+y)$ zeruje się na otoczeniu $K$,
stąd wartość działania $T$ jest zerowa. Zatem obraz
jest funkcją testową.

Dla ustalonego $L\Subset\R^d$, $m\in\N_0$ i
$\varphi\in\mathcal D_L$ oszacowanie z
Lematu~\ref{lem:compact-action21}, zastosowane do
przesunięć funkcji $\varphi$, daje
\[
 p_m(A_T\varphi)\le C_{T,L,m}\,p_{m+r}(\varphi),
 \qquad \supp(A_T\varphi)\subseteq L-K,
\]
gdzie $r$ jest rzędem oszacowania z Lematu~\ref{lem:compact-action21}.
Istotnie, dla $|\beta|\le m$ stosujemy to oszacowanie do
$f(x)=\partial^\beta\varphi(x+y)$; każda pochodna $f$
rzędu co najwyżej $r$ ma moduł ograniczony przez
$p_{m+r}(\varphi)$ niezależnie od $y$.
Można więc wziąć stałą z tego lematu niezależną od $L,m$.
Zbiór $L-K$ jest zwarty jako obraz $L\times K$ przez
ciągłe odwzorowanie $(u,v)\mapsto u-v$; dla $K=\varnothing$
operator jest zerowy.
To dowodzi ciągłości na każdym etapie $\mathcal D_L$,
a z własności uniwersalnej granicy induktywnej
wynika ciągłość na $\mathcal D$. Ciągłe odwzorowanie liniowe
przenosi zbiory ograniczone na ograniczone: dla otoczenia
zera $U$ w przestrzeni docelowej $A_T^{-1}(U)$ jest
otoczeniem zera, a z $B\subseteq t A_T^{-1}(U)$
wynika $A_T(B)\subseteq tU$.
\end{proof}

\begin{definicja}[Splot dystrybucji]
\label{def:convolution21}
Niech $T\in\mathcal E'(\R^d)$ i
$S\in\mathcal D'(\R^d)$. Ich \emph{splotem} nazywamy
dystrybucję $T*S\in\mathcal D'(\R^d)$ określoną wzorem
\[
 \langle T*S,\varphi\rangle
 =\langle S_y,\langle T_x,\varphi(x+y)\rangle\rangle
 =\langle S,A_T\varphi\rangle,
 \qquad\varphi\in\mathcal D(\R^d).
\]
Dla $S\in\mathcal D'(\R^d)$ i $T\in\mathcal E'(\R^d)$
definiujemy również $S*T:=T*S$, jeżeli $S$ nie ma
zwartego nośnika. Gdy oba czynniki mają zwarty nośnik,
zgodność obu możliwych kolejności zostanie dowiedziona
poniżej. W szczególności nie definiujemy tutaj splotu
dowolnych dwóch dystrybucji bez warunków na nośniki.
\end{definicja}

\begin{twierdzenie}[Istnienie, ciągłość i nośnik splotu]
\label{thm:convolution-properties21}
Niech $T\in\mathcal E'(\R^d)$ i
$S\in\mathcal D'(\R^d)$. Wówczas:
\begin{enumerate}
\item splot $T*S$ jest dobrze określoną dystrybucją;
\item $\supp(T*S)\subseteq\supp T+\supp S$;
\item dla ustalonego $T$ odwzorowanie
$S\mapsto T*S$ jest ciągłym operatorem liniowym na
$\mathcal D'(\R^d)$ zarówno dla topologii słabej-*, jak
i dla topologii silnej;
\item jeżeli także $S$ ma zwarty nośnik, to $T*S$ ma
zwarty nośnik.
\end{enumerate}
\end{twierdzenie}
\begin{proof}
Pierwszy punkt jest konsekwencją ciągłości $A_T$:
złożenie ciągłego funkcjonału $S$ z $A_T$ jest ciągłym
funkcjonałem na $\mathcal D$. Oznaczmy $K=\supp T$, $F=\supp S$. Suma $K+F$
jest domknięta: jeżeli $k_n+f_n\to z$, $k_n\in K$,
$f_n\in F$, to z ciągu $(k_n)$ wybieramy podciąg zbieżny
do $k\in K$, a odpowiedni podciąg $(f_n)$ zbiega do
$z-k\in F$. Dla pustego składnika suma jest pusta.
Jeżeli $\supp\varphi$ jest rozłączny z $K+F$, to
$(\supp\varphi-K)\cap F=\varnothing$, gdyż
$y=u-k\in F$ oznaczałoby $u=k+y\in K+F$.
Z Lematu~\ref{lem:convolution-transpose21} wynika więc
$\supp(A_T\varphi)\cap F=\varnothing$, a z
Lematu~\ref{lem:distribution-local19} otrzymujemy
$\langle S,A_T\varphi\rangle=0$.
Stosując to do wszystkich funkcji testowych o nośniku
w otwartym dopełnieniu $K+F$, dostajemy
$\supp(T*S)\subseteq K+F$.

Dla ustalonej funkcji testowej $\varphi$ zachodzi
$|\langle T*S,\varphi\rangle|=
|\langle S,A_T\varphi\rangle|$, co zapewnia
ciągłość słabą-*. Dla zbioru ograniczonego
$B\subseteq\mathcal D$ mamy
\[
 \sup_{\varphi\in B}|\langle T*S,\varphi\rangle|
 =\sup_{\psi\in A_T(B)}|\langle S,\psi\rangle|.
\]
Lemat~\ref{lem:convolution-transpose21} daje
ograniczoność $A_T(B)$, a więc ciągłość silną.
Jeśli oba nośniki są zwarte, ich suma jest zwarta,
co daje ostatni punkt.
\end{proof}

\subsection{Przemienność, łączność i różniczkowanie}

\begin{twierdzenie}[Przemienność i łączność splotu]
\label{thm:convolution-algebra21}
Niech $T,S,U\in\mathcal D'(\R^d)$.
Jeśli co najmniej jeden z czynników $T,S$ ma zwarty nośnik,
to $T*S=S*T$, a wynik nie zależy od sposobu użycia
Definicji~\ref{def:convolution21}. Jeśli co najmniej dwa
spośród $T,S,U$ mają zwarty nośnik, to
\[
 (T*S)*U=T*(S*U).
\]
W szczególności $\mathcal E'(\R^d)$ jest przemienną
algebrą z jedynką $\delta_0$ względem splotu,
a $\mathcal D'(\R^d)$ jest modułem nad tą algebrą.
\end{twierdzenie}
\begin{proof}
Jeżeli $T$ i $S$ mają zwarty nośnik, wybieramy funkcje
odcięcia $\chi_T,\chi_S$, równe $1$ na otoczeniach
odpowiednich nośników. Dla $\varphi\in\mathcal D$ funkcja
$\Phi(x,y)=\chi_T(x)\chi_S(y)\varphi(x+y)$
należy do $\mathcal D(\R^{2d})$. Dla każdego $y$
mnożenie przez $\chi_T$ nie zmienia działania $T_x$
na $\varphi(x+y)$, na mocy Lematu~\ref{lem:compact-action21}.
Zatem $\langle T_x,\Phi(x,y)\rangle=\chi_S(y)A_T\varphi(y)$.
Ponieważ $A_T\varphi$ jest funkcją testową, mnożenie jej
przez $\chi_S$ nie zmienia działania $S$. Otrzymujemy
$\langle S_y,\langle T_x,\Phi(x,y)\rangle\rangle
=\langle T*S,\varphi\rangle$. Analogicznie w odwrotnej
kolejności otrzymujemy $\langle S*T,\varphi\rangle$.
Równość wynika z Lematu~\ref{lem:iterated-distributions21}.
Jeśli tylko jeden czynnik ma zwarty nośnik,
przemienność jest częścią definicji, a powyższy
argument zapewnia jej niesprzeczność.

Załóżmy najpierw, że $T,S\in\mathcal E'(\R^d)$,
a $U\in\mathcal D'(\R^d)$. Z Twierdzenia~
\ref{thm:convolution-properties21} splot $T*S$ ma zwarty
nośnik, więc wszystkie potrzebne sploty są określone.
Dla $\varphi\in\mathcal D(\R^d)$ i ustalonego $z\in\R^d$
funkcja $\varphi_z(w)=\varphi(w+z)$ jest testowa.
Z definicji splotu mamy
\[
\begin{aligned}
 (A_{T*S}\varphi)(z)
 &=\langle T*S,\varphi_z\rangle\\
 &=\langle S_y,\langle T_x,\varphi(x+y+z)\rangle\rangle\\
 &=\langle S_y,(A_T\varphi)(y+z)\rangle
 =(A_S A_T\varphi)(z).
\end{aligned}
\]
Wszystkie działania są dozwolone na mocy
Lematu~\ref{lem:convolution-transpose21}: wewnętrzne
działanie daje przesunięcie funkcji testowej $A_T\varphi$.
Zatem $A_{T*S}=A_S A_T$. Z przemienności dla zwartych
nośników mamy również $A_S A_T=A_T A_S$.
W konsekwencji
\[
\begin{aligned}
 \langle (T*S)*U,\varphi\rangle
 &=\langle U,A_{T*S}\varphi\rangle
 =\langle U,A_S A_T\varphi\rangle\\
 &=\langle S*U,A_T\varphi\rangle
 =\langle T*(S*U),\varphi\rangle .
\end{aligned}
\]
Dla przypadku, w którym $T,U$ mają zwarte nośniki,
przemienność i wykazana już tożsamość dają
\[
\begin{aligned}
 (T*S)*U&=U*(T*S)=(U*T)*S\\
 &=(T*U)*S=T*(U*S)=T*(S*U).
\end{aligned}
\]
Gdy $S,U$ mają zwarte nośniki, analogicznie
\[
 (T*S)*U=U*(S*T)=(U*S)*T=T*(S*U).
\]
W każdym z tych wzorów przynajmniej jeden czynnik
każdego splotu ma zwarty nośnik, gdyż splot dwóch
dystrybucji o zwartych nośnikach ma zwarty nośnik.
W ten sposób wykazaliśmy wszystkie przypadki łączności.

Dla każdej $\varphi\in\mathcal D$ zachodzi
$A_{\delta_0}\varphi=\varphi$, więc $\delta_0*T=T$.
Liniowość w drugim czynniku wynika z liniowości
parowania, a w pierwszym z liniowości
$T\mapsto A_T$ na $\mathcal E'(\R^d)$: dla kombinacji
liniowej dwóch dystrybucji o zwartych nośnikach można
użyć jednego odcięcia równego $1$ w otoczeniu unii
ich nośników.
Stąd przemienna algebra z jedynką oraz struktura
modułu, zgodnie z udowodnioną łącznością.
\end{proof}

\begin{definicja}[Translacja dystrybucji]
\label{def:translation21}
Dla $a\in\R^d$ i $T\in\mathcal D'(\R^d)$ definiujemy
$\tau_aT$ wzorem
$\langle\tau_aT,\varphi\rangle=
\langle T,\varphi(\cdot+a)\rangle$.
Odwzorowanie $\varphi\mapsto\varphi(\cdot+a)$
przekształca $\mathcal D_K$ w $\mathcal D_{K-a}$ i
zachowuje wszystkie seminormy $p_m$. Jest więc ciągłe
na $\mathcal D$ z własności uniwersalnej granicy
induktywnej; $\tau_aT$ jest dobrze określoną dystrybucją.
Dla $f\in L^1_{\mathrm{loc}}(\R^d;\K)$ podstawienie
$z=x+a$ w całce daje
$\langle\tau_aT_f,\varphi\rangle
=\int_{\R^d}f(z-a)\varphi(z)\,\mathrm d\lambda_d(z)$.
Zatem $\tau_aT_f=T_{f(\cdot-a)}$.
\end{definicja}

\begin{przyklad}[Dystrybucje Diraca i przesunięcia]
\label{ex:delta-convolution21}
Dla $a,b\in\R^d$ i $T\in\mathcal D'$ mamy
$\delta_a*T=\tau_aT$, $\delta_a*\delta_b=\delta_{a+b}$.
Dla wielowskaźników $\alpha,\beta\in\N_0^d$ otrzymujemy
$\partial^\alpha\delta_0*\partial^\beta\delta_0
=\partial^{\alpha+\beta}\delta_0$.
Istotnie, $A_{\delta_a}\varphi(y)=\varphi(a+y)$,
co daje pierwszą tożsamość, a dla $T=\delta_b$
także drugą. Dla pochodnych delty mamy
$A_{\partial^\alpha\delta_0}\varphi(y)
=(-1)^{|\alpha|}\partial^\alpha\varphi(y)$.
Zastosowanie $\partial^\beta\delta_0$ do tej funkcji
daje $(-1)^{|\alpha|+|\beta|}
\partial^{\alpha+\beta}\varphi(0)$, czyli działanie
$\partial^{\alpha+\beta}\delta_0$.
\end{przyklad}

\begin{twierdzenie}[Pochodne splotu]
\label{thm:convolution-derivative21}
Niech $T,S\in\mathcal D'(\R^d)$ i niech co najmniej
jeden z nich ma zwarty nośnik. Dla każdego
$\alpha\in\N_0^d$ zachodzi
\[
 \partial^\alpha(T*S)
 =(\partial^\alpha T)*S
 =T*(\partial^\alpha S).
\]
W szczególności dla wielomianu $P$ w $d$ zmiennych
i operatora stałowspółczynnikowego $P(\partial)$
operator splotu z ustaloną dystrybucją o zwartym
nośniku komutuje z $P(\partial)$.
\end{twierdzenie}
\begin{proof}
Załóżmy, że $T$ ma zwarty nośnik $K$. Ustalmy
$j\in\{1,\ldots,d\}$.
Z Twierdzenia~\ref{thm:dist-deriv-cont20} wszystkie
pochodne $T$ mają nośniki zawarte w $K$, więc ich sploty
z $S$ są określone. Dla $f\in C^\infty(\R^d)$
wzór na pochodną dystrybucji pozostaje prawdziwy
dla przedłużonego działania:
\[
 \langle\partial_jT,f\rangle
 =-\langle T,\partial_jf\rangle.
\]
Istotnie, bierzemy $\chi\in\mathcal D$ równą $1$
w otoczeniu $K$. Wówczas
$\langle\partial_jT,\chi f\rangle
=-\langle T,\chi\partial_jf\rangle
-\langle T,(\partial_j\chi)f\rangle$.
Ostatni składnik znika, gdyż nośnik funkcji
$(\partial_j\chi)f$ jest rozłączny z $K$.
Pierwszy składnik jest przedłużonym działaniem
$-T$ na $\partial_jf$. Dla
$\varphi\in\mathcal D$ i $j\in\{1,\ldots,d\}$ mamy
\[
 \begin{aligned}
 \langle\partial_j(T*S),\varphi\rangle
 &=-\langle S_y,\langle T_x,\partial_j\varphi(x+y)\rangle\rangle\\
 &=\langle S_y,\langle(\partial_jT)_x,\varphi(x+y)\rangle\rangle
 =\langle(\partial_jT)*S,\varphi\rangle.
 \end{aligned}
\]
W drugim kroku różniczkowanie dotyczy zmiennej $x$.
Aby uzyskać $T*(\partial_jS)$, zauważamy, że
$\partial_{y_j}A_T\varphi=A_T(\partial_j\varphi)$,
więc definicja pochodnej $S$ daje tę samą wartość.
Dla $\alpha=0$ wzór jest oczywisty. Indukcja względem
$|\alpha|$, wykorzystująca zachowanie zwartości nośnika
pochodnych $T$, daje obie tożsamości dla każdego
wielowskaźnika. Gdy zwarty nośnik ma tylko $S$,
stosujemy przemienność. Dla wielomianu
$P(\xi)=\sum_\alpha c_\alpha\xi^\alpha$ sumujemy
równości $T*(\partial^\alpha S)=\partial^\alpha(T*S)$
z wagami $c_\alpha$, co daje ostatnie stwierdzenie.
\end{proof}

\begin{przyklad}[Inkluzja nośników może być ostra]
\label{ex:support-strict21}
Niech $T=\delta_0-\delta_1\in\mathcal E'(\R)$,
a $S=T_1$ będzie dystrybucją odpowiadającą funkcji
stałej $1$. Wtedy $T*S=T_1-T_1=0$,
mimo że $\supp T+\supp S=\R$. Dlatego
inkluzji nośników nie można zastąpić równością.
\end{przyklad}

\subsection{Splot z funkcją testową i aproksymacja jedności}

\noindent\emph{Powiązanie tematyczne w Bibliotece matematycznej:} \atlasref{PDE-07::thm:pde7-mollify}{PDE-07: wygładzanie przez splot}.


\begin{twierdzenie}[Wygładzanie dystrybucji]
\label{thm:smoothing21}
Niech $T\in\mathcal D'(\R^d)$ i
$\rho\in\mathcal D(\R^d)$. Wówczas splot
$T*\rho$ jest dystrybucją regularną odpowiadającą
funkcji $u\in C^\infty(\R^d)$ danej przez
\[
 u(x)=\langle T_y,\rho(x-y)\rangle.
\]
Dla każdego wielowskaźnika $\alpha$ zachodzi
$\partial^\alpha u(x)
=\langle T_y,(\partial^\alpha\rho)(x-y)\rangle$.
Ponadto $\supp u\subseteq\supp T+\supp\rho$.
\end{twierdzenie}
\begin{proof}
Ustalmy zwarty zbiór $M\subseteq\R^d$ i oznaczmy
$K=\supp\rho$. Dla $x\in M$ funkcje
$y\mapsto\rho(x-y)$ mają nośniki w $M-K$.
Przy badaniu ilorazów różnicowych powiększamy $M$
do zwartego otoczenia badanego punktu. Względem każdej
seminormy $p_m$ iloraz różnicowy w kierunku $e_j$
zbiega do $y\mapsto(\partial_j\rho)(x-y)$:
z twierdzenia Newtona--Leibniza błąd jest ograniczony
przez moduł ciągłości pochodnych $\rho$ rzędu co najwyżej
$m+1$. Wszystkie te funkcje mają wspólny zwarty nośnik.
Ciągłość $T$ na odpowiednim etapie $\mathcal D_L$
pozwala przenieść różniczkowanie pod znak parowania.
Powtarzając argument dla kolejnych pochodnych,
otrzymujemy podany wzór; ciągłość tych pochodnych
wynika z jednostajnej ciągłości pochodnych $\rho$.
Jeśli $x\notin\supp T+K$, to
$(x-K)\cap\supp T=\varnothing$. Lokalność dystrybucji
(Lemat~\ref{lem:distribution-local19}) daje $u(x)=0$.
Ponieważ $\supp T+K$ jest domknięty, otrzymujemy
żądaną inkluzję nośników.

Ustalmy $\varphi\in\mathcal D(\R^d)$ i połóżmy
$G(x,y)=\rho(x-y)\varphi(x)$. Jego nośnik zawiera się
w zwartym zbiorze
\[
 \supp\varphi\times(\supp\varphi-K),
\]
więc $G\in\mathcal D(\R^{2d})$.
Niech $T_1$ oznacza dystrybucję regularną odpowiadającą
funkcji stałej $1$. Stosujemy
Lemat~\ref{lem:iterated-distributions21} do $T_1$ w zmiennej
$x$ i $T$ w zmiennej $y$. Daje to
\[
 \begin{aligned}
 \int_{\R^d}u(x)\varphi(x)\,\mathrm d\lambda_d(x)
 &=\left\langle T_y,
      \int_{\R^d}\rho(x-y)\varphi(x)\,\mathrm d\lambda_d(x)
   \right\rangle\\
 &=\langle T,\check\rho*\varphi\rangle
 =\langle T* T_\rho,\varphi\rangle.
 \end{aligned}
\]
Tutaj $\check\rho(z)=\rho(-z)$, a klasyczny splot
funkcji definiujemy przez
$(\check\rho*\varphi)(y)=\int_{\R^d}\check\rho(y-x)
\varphi(x)\,\mathrm d\lambda_d(x)$.
Jest to dokładnie $A_{T_\rho}\varphi(y)$: podstawienie
$x=y+z$ daje
\[
 A_{T_\rho}\varphi(y)=
 \int_{\R^d}\rho(z)\varphi(y+z)\,\mathrm d\lambda_d(z).
\]
Ostatnia równość wynika więc z definicji splotu
dystrybucji i jego przemienności. Nie wymagamy
zwartego nośnika $T$. W rezultacie $T*T_\rho=T_u$.
\end{proof}

\begin{definicja}[Standardowa aproksymacja jedności]
\label{def:mollifier21}
Wybierzmy $\rho\in\mathcal D(\R^d)$ taką, że
$\rho\ge0$, $\supp\rho\subseteq\overline{B(0,1)}$
oraz $\int_{\R^d}\rho(x)\,\mathrm d\lambda_d(x)=1$.
Istnienie $\rho$ wynika z Lematu~\ref{lem:cutoffs17}:
wybieramy nieujemną funkcję testową w kuli jednostkowej,
równą $1$ na pewnej mniejszej kuli, i dzielimy ją przez
jej dodatnią całkę. Dla $\varepsilon>0$ piszemy
$\rho_\varepsilon(x)=\varepsilon^{-d}\rho(x/\varepsilon)$.
Zmiana zmiennych $x=\varepsilon z$ pokazuje, że
$\int_{\R^d}\rho_\varepsilon(x)\,\mathrm d\lambda_d(x)=1$,
a $\supp\rho_\varepsilon\subseteq\overline{B(0,\varepsilon)}$.
\end{definicja}

\begin{twierdzenie}[Regularyzacja jest silnie zbieżna]
\label{thm:strong-mollification21}
Dla każdego $T\in\mathcal D'(\R^d)$ i każdej
standardowej aproksymacji jedności z
Definicji~\ref{def:mollifier21} dystrybucje
$T*\rho_\varepsilon$ są reprezentowane przez funkcje
$C^\infty(\R^d)$ oraz
\[
 T*\rho_\varepsilon\longrightarrow T
 \quad\text{w topologii silnej }\beta(\mathcal D',\mathcal D)
 \quad(\varepsilon\downarrow0).
\]
Jeśli $T$ ma zwarty nośnik $K$, to
$\supp(T*\rho_\varepsilon)\subseteq K+
\overline{B(0,\varepsilon)}$, więc dla małych
$\varepsilon$ funkcje wygładzające mają nośniki
w jednym zbiorze zwartym.
\end{twierdzenie}
\begin{proof}
Z Twierdzenia~\ref{thm:smoothing21} mamy wygładzanie.
Dla $\varphi\in\mathcal D$ zachodzi
\[
 \langle T*\rho_\varepsilon-T,\varphi\rangle
 =\langle T,\check\rho_\varepsilon*\varphi-\varphi\rangle.
\]
Niech $B\subseteq\mathcal D$ będzie niepustym zbiorem
ograniczonym (dla pustego zbioru warunek jest trywialny).
Z Twierdzenia~\ref{thm:D-bounded18}
wynika istnienie $K\Subset\R^d$ zawierającego nośniki
wszystkich $\varphi\in B$, przy czym
$M_m(B)=\sup_{\varphi\in B}p_m(\varphi)<\infty$
dla wszystkich $m$. Dla $0<\varepsilon\le1$
różnice $\check\rho_\varepsilon*\varphi-\varphi$
mają nośniki w jednym zwartym zbiorze
$L=K+\overline{B(0,1)}$. Dla $m\in\N_0$ i $|\alpha|\le m$ z normalizacji jądra
i podstawienia $z=x/\varepsilon$ mamy
\[
\begin{aligned}
 \partial^\alpha(\check\rho_\varepsilon*\varphi-\varphi)(y)
 =\int_{\R^d}\rho(z)
 \bigl(\partial^\alpha\varphi(y+\varepsilon z)
 -\partial^\alpha\varphi(y)\bigr)\,\mathrm d\lambda_d(z).
\end{aligned}
\]
Różniczkowanie pod całką jest uzasadnione zwartym
nośnikiem $\rho$ i ograniczonością pochodnych $\varphi$.
Z twierdzenia Newtona--Leibniza wzdłuż odcinka
$t\mapsto y+t\varepsilon z$ otrzymujemy
\[
 |\partial^\alpha\varphi(y+\varepsilon z)
 -\partial^\alpha\varphi(y)|
 \le\varepsilon\|z\|_1 p_{m+1}(\varphi),
 \qquad \|z\|_1=\sum_{j=1}^d|z_j|.
\]
Nierówność obowiązuje również dla funkcji zespolonych.
Biorąc supremum po $y$ i $\varphi\in B$ oraz maksimum
po $|\alpha|\le m$, dostajemy
\[
 \sup_{\varphi\in B}p_m(\check\rho_\varepsilon*\varphi
 -\varphi)
 \le C_{\rho,d}\varepsilon\,M_{m+1}(B),
 \quad
 C_{\rho,d}=\int_{\R^d}\|z\|_1\rho(z)\,\mathrm d\lambda_d(z).
\]
Ciągłość $T|_{\mathcal D_L}$ oznacza, że dla pewnych
$m,C$ mamy $|\langle T,\psi\rangle|\le Cp_m(\psi)$
gdy $\psi\in\mathcal D_L$. Dlatego
\[
 \sup_{\varphi\in B}|\langle T*\rho_\varepsilon-T,
 \varphi\rangle|
 \le CC_{\rho,d}\varepsilon M_{m+1}(B)\longrightarrow0.
\]
Ponieważ $B$ jest dowolnym zbiorem ograniczonym,
jest to zbieżność silna. Oszacowanie nośnika wynika
z Twierdzenia~\ref{thm:smoothing21}.
\end{proof}

\begin{wniosek}[Gęstość funkcji testowych dla zwartych nośników]
\label{cor:dense-test-compact21}
Niech $T\in\mathcal E'(\R^d)$. Istnieje ciąg
$(f_n)\subseteq\mathcal D(\R^d)$ taki, że
$T_{f_n}\to T$ silnie w $\mathcal D'(\R^d)$.
Można wziąć $f_n(x)=\langle T_y,\rho_{1/n}(x-y)\rangle$,
czyli funkcję reprezentującą dystrybucję $T*T_{\rho_{1/n}}$.
\end{wniosek}
\begin{proof}
Funkcje $f_n$ są gładkie i mają nośniki zwarte
z poprzedniego twierdzenia. Zbieżność silna także
wynika z tego twierdzenia.
\end{proof}

\begin{przyklad}[Regularyzacja delty Diraca]
\label{ex:delta-smooth21}
Dla $a\in\R^d$ i standardowej aproksymacji jedności
$\rho_\varepsilon$ zachodzi
$\delta_a*\rho_\varepsilon=T_{\rho_\varepsilon(\cdot-a)}$.
Funkcje te stają się coraz bardziej skupione
w otoczeniu $a$, a jako dystrybucje zbiegają
silnie do $\delta_a$. Nie ma tu zbieżności w topologii $L^1_{\mathrm{loc}}$
do żadnej funkcji $f$: gdyby taka zbieżność zachodziła,
to dla każdej $\varphi\in\mathcal D(\R^d)$
\[
 |\langle T_{\rho_\varepsilon(\cdot-a)}-T_f,\varphi\rangle|
 \le p_0(\varphi)
 \int_{\supp\varphi}|\rho_\varepsilon(x-a)-f(x)|
 \,\mathrm d\lambda_d(x)\longrightarrow0.
\]
Jednoznaczność granicy słabej-* dałaby $T_f=\delta_a$,
co jest niemożliwe. Istotnie, ograniczając obie dystrybucje do
$\R^d\setminus\{a\}$ i stosując
Stwierdzenie~\ref{prop:regular-injective19}, otrzymalibyśmy, że
$f=0$ prawie wszędzie poza $\{a\}$, a zatem prawie wszędzie
w $\R^d$, podczas gdy $\delta_a\ne0$.
\end{przyklad}

\subsection{Ograniczenia operacji splotu i powiązania z dalszą teorią}

\begin{uwaga}[Kiedy splot nie jest zdefiniowany]
\label{rem:convolution-limits21}
Warunek zwartego nośnika przynajmniej jednego
czynnika jest wystarczający, ale niekonieczny.
Dla funkcji odpowiednio całkowalnych splot może istnieć
bez tego warunku. Natomiast nie można rozszerzyć splotu na wszystkie pary
dystrybucji tak, aby zachować jego dotychczasowe wartości
i obie równości dotyczące różniczkowania.
Już dla $d=1$ niech $H$ będzie funkcją Heaviside'a,
a $T_1$ dystrybucją funkcji stałej $1$.
Z rozdziału~20 mamy $\partial T_H=\delta_0$ i
$\partial T_1=0$. Gdyby istniało takie rozszerzenie,
to jednocześnie
\[
 \partial(T_1*T_H)=(\partial T_1)*T_H=0
 \quad\text{oraz}\quad
 \partial(T_1*T_H)=T_1*\delta_0=T_1,
\]
co jest sprzeczne z $T_1\ne0$.
Dla $d>1$ zamiast $T_H$ bierzemy funkcjonał $V$
określony dla $\varphi\in\mathcal D(\R^d)$ przez
$\langle V,\varphi\rangle=\int_0^\infty
\varphi(t,0,\ldots,0)\,\mathrm d\lambda_1(t)$.
Na każdym etapie $\mathcal D_K$ działanie to jest
ograniczone przez $p_0(\varphi)$ razy długość
przedziału zawierającego rzut $K$ na pierwszą oś,
więc jest ciągłe na każdym etapie. Liniowość wynika z
liniowości całki, a własność uniwersalna granicy induktywnej
daje $V\in\mathcal D'(\R^d)$. Dla $\varphi\in\mathcal D(\R^d)$
z twierdzenia Newtona--Leibniza mamy
\[
 \langle\partial_1V,\varphi\rangle
 =-\int_0^\infty(\partial_1\varphi)(t,0,\ldots,0)
 \,\mathrm d\lambda_1(t)=\varphi(0).
\]
Zatem $\partial_1 V=\delta_0$, podczas gdy
$\partial_1 T_1=0$. Uzyskujemy tę samą sprzeczność.

Na przykład formalny splot dwóch funkcji stałych
na $\R^d$ wymagałby całkowania stałej $1$
po całej $\R^d$, co nie daje skończonej wartości.
Dalsze rozszerzenia opierają się między innymi
na warunkach dotyczących nośników, zapewniających
właściwość odwzorowania dodawania na ich iloczynie
(tzn. zwartość przeciwobrazu każdego zbioru zwartego),
lub na dodatkowych własnościach analitycznych czynników.
\end{uwaga}

\begin{uwaga}[Dlaczego jest to zagadnienie analizy funkcjonalnej]
\label{rem:convolution-tvs21}
Kluczowe kroki nie polegały na rozwiązywaniu równań
różniczkowych. Użyliśmy ciągłości operatora
$A_T$ na przestrzeni funkcji testowych, transpozycji
na jej przestrzeń dualną, przenoszenia zbiorów ograniczonych
oraz zbieżności silnej w tej przestrzeni dualnej.
Są to konstrukcje z teorii przestrzeni liniowo-topologicznych.
Dalsze zastosowania splotu do rozwiązywania równań
różniczkowych pozostawiamy odrębnemu skryptowi.
\end{uwaga}

% Materiał z AF 25; fragment do rozdziału o splocie.

\subsection{Aproksymacja w przestrzeniach Lebesgue'a}

\begin{twierdzenie}[Gęstość funkcji testowych]
\label{thm:dt-test-density}
Dla $1\le p<\infty$ funkcje $C_c^\infty(\R^d)$
są gęste w $L^p(\R^d)$.
\end{twierdzenie}
\begin{proof}
Obcięcie funkcji do kuli i do ograniczonych wartości
zbiega do niej w $L^p$ przez zbieżność zmajoryzowaną.
Na zbiorze skończonej miary ograniczoną funkcję
można aproksymować jednostajnie funkcjami prostymi
(oddzielnie część rzeczywistą i urojoną).
Dla mierzalnego ograniczonego zbioru $E$ regularność
miary Lebesgue'a daje zwarty $F\subset E$ oraz
ograniczony otwarty $O\supset E$ z
$\lambda_d(O\setminus F)<\eta$ dla dowolnego $\eta>0$.
Gdy $F\ne\varnothing$, wybieramy $r>0$ z
$2r<\operatorname{dist}(F,\R^d\setminus O)$.
Niech $\rho$ będzie nieujemną gładką funkcją o nośniku
w $B(0,r/2)$ i całce $1$. Funkcja
$\psi=\chi_{\{x:\operatorname{dist}(x,F)<r\}}*\rho$
jest gładka na mocy twierdzenia \ref{thm:smoothing21},
ma wartości w $[0,1]$, jest równa $1$ na $F$
i ma zwarty nośnik w $O$. Gdy $F=\varnothing$,
bierzemy $\psi=0$. W obu przypadkach
$\|\chi_E-\psi\|_p^p\le\lambda_d(O\setminus F)<\eta$.
Aproksymując każdy składnik funkcji prostej i używając
nierówności trójkąta, otrzymujemy tezę.
\end{proof}

\begin{twierdzenie}[Aproksymacja jedności w $L^p$]
\label{thm:dt-approx-lp}
Niech $1\le p<\infty$, $f\in L^p(\R^d)$ i niech
$\rho_\varepsilon$ będzie standardową aproksymacją
jedności z definicji \ref{def:mollifier21}.
Wtedy $\|f*\rho_\varepsilon-f\|_p\to0$.
Jeżeli $f$ ma zwarty nośnik, to $f*\rho_\varepsilon$
jest funkcją testową.
\end{twierdzenie}
\begin{proof}
Najpierw dla $v\in C_c^\infty$ mamy
$\|v(\cdot-y)-v\|_p\to0$ przy $y\to0$:
dla $|y|\le1$ różnica ma nośnik w jednym zbiorze
zwartym, a jednostajna ciągłość daje zbieżność
jednostajną. Dla dowolnego $f\in L^p$
i takiej funkcji $v$ nierówność trójkąta
oraz niezmienniczość normy względem translacji dają
\[
 \|f(\cdot-y)-f\|_p\le2\|f-v\|_p+\|v(\cdot-y)-v\|_p.
\]
Gęstość z twierdzenia \ref{thm:dt-test-density}
dowodzi ciągłości translacji w normie $L^p$.
Nierówność Jensena dla probabilistycznej miary
$\rho_\varepsilon(y)\,\mathrm d\lambda_d(y)$
i twierdzenie Tonellego dają
\[
 \|f*\rho_\varepsilon-f\|_p^p
 \le \int_{\R^d}\rho_\varepsilon(y)
            \|f(\cdot-y)-f\|_p^p\,\mathrm d\lambda_d(y)
 \le\sup_{|y|\le\varepsilon}\|f(\cdot-y)-f\|_p^p.
\]
Lewa strona jest dobrze określona: zastosowanie
Jensena i Tonellego do modułu $f(\cdot-y)$
pokazuje, że splot istnieje prawie wszędzie
i należy do $L^p$. Ostatnie supremum dąży do zera
przez ciągłość translacji w zerze.
W przypadku zwartego nośnika gładkość i zwartość
nośnika splotu wynikają z twierdzenia
\ref{thm:smoothing21}.
\end{proof}

